\documentclass{article}
\usepackage{booktabs}
\usepackage{tabularx}
\usepackage[margin=1in]{geometry}
\begin{document}

\begin{table}[tbp] \centering
\newcolumntype{R}{>{\raggedleft\arraybackslash}X}
\newcolumntype{L}{>{\raggedright\arraybackslash}X}
\newcolumntype{C}{>{\centering\arraybackslash}X}

\caption{__000006}
\label{__00000A}
\begin{tabularx}{\linewidth}{@{}lCCCCCC@{}}

\toprule
&{(1)}&{(2)}&{(3)}&{(4)}&{(5)}&{(6)} \tabularnewline \midrule
{Vote no.}&{Year}&{Title (short)}&{Vineyard coef}&{(SE)}&{p}&{} \tabularnewline
\midrule \addlinespace[\belowrulesep]
56&1900&Gesetz zur Kranken-, Unfall- und Militärversicherung&--203.6&(312)&0.522& \tabularnewline
57&1900&Initiative «für die Proporzwahl des Nationalrates»&--478.9&(304)&0.130& \tabularnewline
58&1900&Initiative für die Volkswahl des Bundesrates&--431.1&(346)&0.226& \tabularnewline
59&1902&Unterstützung der Primarschule durch den Bund&383.4&(482)&0.435& \tabularnewline
60&1903&Zolltarifgesetz&737.3&(737)&0.328& \tabularnewline
61&1903&Bundesstrafrecht (Anstiftung Militärpflichtiger zu Ver&99.6&(314)&0.754& \tabularnewline
62&1903&Initiative «für die Wahl des Nationalrates aufgrund d&--3.3&(655)&0.996& \tabularnewline
63&1903&Artikel über die Regulierung des Alkoholhandels&--51.7&(356)&0.886& \tabularnewline
64&1905&Ausdehnung des Erfinderschutzes&352.7&(387)&0.373& \tabularnewline
65&1906&Lebensmittelgesetz&1285.6**&(604)&0.045& \tabularnewline
66&1907&Militärorganisation der schweizerischen Eidgenossensch&300.6&(355)&0.407& \tabularnewline
67&1908&Verfassungsartikel zur Gesetzgebung über das Gewerbewe&76.0&(445)&0.866& \tabularnewline
68&1908&Initiative «für ein Absinthverbot»&484.4**&(200)&0.024&TREAT \tabularnewline
69&1908&Verfassungsartikel über die Wasserkraft und Elektrizit&65.4&(205)&0.753& \tabularnewline
70&1910&Initiative «für die Proporzwahl des Nationalrates»&--61.5&(524)&0.908& \tabularnewline
\bottomrule \addlinespace[\belowrulesep]

\end{tabularx}
\\ \parbox{\linewidth}{\footnotesize Notes: KEY-spec OLS (yes\\_pct on vineyard\\_per\\_cap + french\\_share + catholic\\_share, HC3 SEs) run separately on each of the 15 federal popular votes between 1900 and 1910. The treatment vote (\\#68, 1908 absinthe ban) is marked TREAT. Falsification logic: if vineyard\\_per\\_cap predicts yes-vote shares broadly, the absinthe finding is spurious; if only \\#68 plus substantively related votes show non-null coefficients, the wine-protection mechanism is issue-specific. Vote \\#65 (Lebensmittelgesetz, 1906) is NOT a clean placebo: this Federal Act established the alcohol-regulation authority later invoked against absinthe and was supported by wine producers because it cracked down on wine adulteration and substitute beverages. Treat \\#65 as the regulatory prequel to \\#68, not an independent comparison. Vote \\#63 (1903 alcohol-trade regulation, distinct earlier coalition that failed) is the cleaner alcohol-regulation null. Stars: * p<0.10, ** p<0.05, *** p<0.01.}
\end{table}
\end{document}
